\documentclass[12pt, a4paper]{article}
\usepackage{../../../myconfig} % Gọi file cấu hình từ ngoài gốc vào

% ==========================================
% BẮT ĐẦU NỘI DUNG TÀI LIỆU
% ==========================================
\begin{document}

% Truyền 3 tham số: {Nhãn cam} {Chữ to chính giữa} {Chữ ở góc phải Header}
\titlebox{Chủ đề: Tích phân}{Tích phân}{Nguyên hàm - Tích phân: Tích phân}

\drawdashlines{20}
\newpage
\drawdashlines{25}
\newpage
\drawdashlines{25}
\newpage
\drawdashlines{25}
\newpage

% --- CÂU 1 ---
\questionLinesManual{6}{1}{Nếu \( \displaystyle \int_{1}^{2} f(x) \, dx = -2 \) và \( \displaystyle \int_{2}^{3} f(x) \, dx = 1 \) thì \( \displaystyle \int_{1}^{3} f(x) \, dx \) bằng}{
    \begin{tasks}(4)
        \task $-3$
        \task $-1$
        \task $1$
        \task $3$
    \end{tasks}
}
% --- CÂU 2 ---
\questionLinesManual{6}{2}{Nếu \( \displaystyle \int_{0}^{1} f(x) \, dx = 4 \) thì \( \displaystyle \int_{0}^{1} 2f(x) \, dx \) bằng}{
    \begin{tasks}(4)
        \task $16$
        \task $4$
        \task $2$
        \task $8$
    \end{tasks}
}

% --- CÂU 3 ---
\questionLinesManual{8}{3}{Biết \( F(x) = x^2 \) là một nguyên hàm của hàm số \( f(x) \) trên \( \mathbb{R} \). Giá trị của \( \displaystyle \int_{1}^{2} [2 + f(x)] \, dx \) bằng}{
    \begin{tasks}(4)
        \task $5$
        \task $3$
        \task $\dfrac{13}{3}$
        \task $\dfrac{7}{3}$
    \end{tasks}
}

% --- CÂU 4 ---
\questionLinesManual{6}{4}{Biết \( \displaystyle \int_{2}^{3} f(x) \, dx = 4 \) và \( \displaystyle \int_{2}^{3} g(x) \, dx = 1 \). Khi đó \( \displaystyle \int_{2}^{3} [f(x) - g(x)] \, dx \) bằng:}{
    \begin{tasks}(4)
        \task $-3$
        \task $3$
        \task $4$
        \task $5$
    \end{tasks}
}

% --- CÂU 5 ---
\questionLinesManual{8}{5}{Biết \( \displaystyle \int_{0}^{1} [f(x) + 2x] \, dx = 2 \). Khi đó \( \displaystyle \int_{0}^{1} f(x) \, dx \) bằng:}{
    \begin{tasks}(4)
        \task $1$
        \task $4$
        \task $2$
        \task $0$
    \end{tasks}
}

% --- CÂU 6 ---
\questionLinesManual{7}{6}{Cho \( \displaystyle \int_{-2}^{2} f(x) \, dx = 1 \), \( \displaystyle \int_{-2}^{4} f(t) \, dt = -4 \). Tính \( \displaystyle \int_{2}^{4} f(y) \, dy \).}{
    \begin{tasks}(4)
        \task $I = 5$
        \task $I = -3$
        \task $I = 3$
        \task $I = -5$
    \end{tasks}
}

% --- CÂU 7 ---
\questionLinesManual{6}{7}{Cho \( \displaystyle \int_{0}^{1} f(x) \, dx = -1 \); \( \displaystyle \int_{0}^{3} f(x) \, dx = 5 \). Tính \( \displaystyle \int_{1}^{3} f(x) \, dx \).}{
    \begin{tasks}(4)
        \task $1$
        \task $4$
        \task $6$
        \task $5$
    \end{tasks}
}

% --- CÂU 8 ---
\questionLinesManual{6}{8}{Cho \( \displaystyle \int_{1}^{2} f(x) \, dx = -3 \) và \( \displaystyle \int_{2}^{3} f(x) \, dx = 4 \). Khi đó \( \displaystyle \int_{1}^{3} f(x) \, dx \) bằng}{
    \begin{tasks}(4)
        \task $12$
        \task $7$
        \task $1$
        \task $-12$
    \end{tasks}
}

% --- CÂU 9 ---
\questionLinesManual{10}{9}{Nếu \( F'(x) = \dfrac{1}{2x-1} \) và \( F(1) = 1 \) thì giá trị của \( F(4) \) bằng}{
    \begin{tasks}(4)
        \task $\ln 7$
        \task $1 + \dfrac{1}{2}\ln 7$
        \task $\ln 3$
        \task $1 + \ln 7$
    \end{tasks}
}

% --- CÂU 10 ---
\questionLinesManual{6}{10}{Tính tích phân \( I = \displaystyle \int_{0}^{2} (2x + 1) \, dx \).}{
    \begin{tasks}(4)
        \task $I = 5$
        \task $I = 6$
        \task $I = 2$
        \task $I = 4$
    \end{tasks}
}

% --- CÂU 11 ---
\questionLinesManual{6}{11}{Với \( a, b \) là các tham số thực. Giá trị tích phân \( \displaystyle \int_{0}^{b} (3x^2 - 2ax - 1) \, dx \) bằng}{
    \begin{tasks}(4)
        \task \( b^3 - b^2a - b \)
        \task \( b^3 + b^2a + b \)
        \task \( b^3 - ba^2 - b \)
        \task \( 3b^2 - 2ab - 1 \)
    \end{tasks}
}

% --- CÂU 12 ---
\questionLinesManual{10}{12}{Cho hàm số \( f(x) \). Biết \( f(0) = 4 \) và \( f'(x) = 2\sin x + 1, \, \forall x \in \mathbb{R} \), khi đó \( \displaystyle \int_{0}^{\frac{\pi}{4}} f(x) \, dx \) bằng}


% --- CÂU 13 ---
\questionLinesManual{10}{13}{Biết \( \displaystyle \int_{1}^{3} \dfrac{x+2}{x} \, dx = a + b\ln c \), với \( a, b, c \in \mathbb{Z}, c < 9 \). Tính tổng \( S = a + b + c \).}{
    \begin{tasks}(4)
        \task \( S = 7 \)
        \task \( S = 5 \)
        \task \( S = 8 \)
        \task \( S = 6 \)
    \end{tasks}
}

% --- CÂU 14 ---
\questionLinesManual{12}{14}{Biết \( I = \displaystyle \int_{-1}^{0} \dfrac{3x^2 + 5x - 1}{x - 2} \, dx = a\ln \dfrac{2}{3} + b, \, (a, b \in \mathbb{R}) \). Khi đó giá trị của \( a + 4b \) bằng}{
    \begin{tasks}(4)
        \task $50$
        \task $60$
        \task $59$
        \task $40$
    \end{tasks}
}


% --- CÂU 15 ---
\questionLinesManual{23}{15}{Cho \( \displaystyle \int_{0}^{1} \left( \dfrac{1}{x+1} - \dfrac{1}{x+2} \right) dx = a\ln 2 + b\ln 3 \) với \( a, b \) là các số nguyên. Mệnh đề nào dưới đây đúng?}{
    \begin{tasks}(4)
        \task \( a + 2b = 0 \)
        \task \( a + b = 2 \)
        \task \( a - 2b = 0 \)
        \task \( a + b = -2 \)
    \end{tasks}
}

% --- CÂU 16 ---
\questionLinesManual{23}{16}{Biết \( \displaystyle \int_{1}^{2} \dfrac{dx}{(x+1)(2x+1)} = a\ln 2 + b\ln 3 + c\ln 5 \). Khi đó giá trị \( a + b + c \) bằng}{
    \begin{tasks}(4)
        \task $-3$
        \task $2$
        \task $1$
        \task $0$
    \end{tasks}
}

% --- CÂU 17 ---
\questionLinesManual{23}{17}{Cho \( \displaystyle \int_{1}^{3} \dfrac{x+3}{x^2 + 3x + 2} \, dx = a\ln 2 + b\ln 3 + c\ln 5 \), với \( a, b, c \) là các số nguyên. Giá trị của \( a+b+c \) bằng}{
    \begin{tasks}(4)
        \task $0$
        \task $2$
        \task $3$
        \task $1$
    \end{tasks}
}

% --- CÂU 18 ---
\questionLinesManual{23}{18}{Cho \( \displaystyle \int_{3}^{4} \dfrac{5x-8}{x^2 - 3x + 2} \, dx = a\ln 3 + b\ln 2 + c\ln 5 \), với \( a, b, c \) là các số hữu tỉ. Giá trị của \( 2^{a-3b+c} \) bằng}{
    \begin{tasks}(4)
        \task $12$
        \task $6$
        \task $1$
        \task $64$
    \end{tasks}
}








\end{document}
